\documentclass[12pt]{article}
\usepackage[utf8]{inputenc}
\usepackage[T1]{fontenc}
\usepackage[margin=1in]{geometry}
\usepackage{setspace}
\usepackage{xcolor}
\usepackage[colorlinks=true, urlcolor=blue!60!black, linkcolor=blue!60!black]{hyperref}
\setlength{\parindent}{2em}
\setlength{\parskip}{1.25em}
\renewcommand{\baselinestretch}{1.0}
\begin{document}
\begin{center}
{\LARGE \textit{Washington v. Glucksberg}, 521 U.S. 702 (1997)\par}
\vspace{0.8em}
{\large Case Brief\par}
\end{center}
\vspace{1em}

\section{Memory Jogger}

Because an asserted right to physician-assisted suicide is not ``deeply rooted in this Nation's history and tradition,'' Washington's criminal ban on aiding a suicide does not violate the Fourteenth Amendment's Due Process Clause (\href{https://www.courtlistener.com/opinion/118144/washington-v-glucksberg/}{CourtListener}; \href{https://supreme.justia.com/cases/federal/us/521/702/}{Justia}).

\section{Detailed Case Facts}

Since 1854, when its first Territorial Legislature outlawed ``assisting another in the commission of self-murder,'' Washington had criminalized assisting a suicide (\href{https://supreme.justia.com/cases/federal/us/521/702/}{Justia}). Under the 1975 criminal code revision at issue, Wash. Rev. Code \S~9A.36.060(1) (1994), ``[a] person is guilty of promoting a suicide attempt when he knowingly causes or aids another person to attempt suicide,'' a felony punishable by up to five years' imprisonment and a fine of up to \$10,000 (\href{https://supreme.justia.com/cases/federal/us/521/702/}{Justia}; \href{https://www.law.cornell.edu/supct/html/96-110.ZO.html}{Cornell LII}). By contrast, Washington's Natural Death Act of 1979 expressly recognized an adult patient's right to direct the withholding or withdrawal of life-sustaining treatment in a terminal or permanent unconscious condition, declaring that such withdrawal ``shall not, for any purpose, constitute a suicide,'' Wash. Rev. Code \S~70.122.070(1), while excluding mercy killing and---following Washington voters' 1991 rejection of ballot Initiative 119---expressly excluding physician-assisted suicide, Wash. Rev. Code \S~70.122.100 (\href{https://supreme.justia.com/cases/federal/us/521/702/}{Justia}).

In January 1994, four Washington physicians who occasionally treated terminally ill, suffering patients---Dr. Harold Glucksberg, Dr. Abigail Halperin, Dr. Thomas A. Preston, and Dr. Peter Shalit---joined with three gravely ill, pseudonymous patients (Jane Roe, John Doe, and James Poe, all of whom died while the litigation was pending) and Compassion in Dying, a nonprofit organization that counseled individuals considering physician-assisted suicide, to sue the State of Washington and its Attorney General (\href{https://supreme.justia.com/cases/federal/us/521/702/}{Justia}; \href{https://www.law.cornell.edu/supct/html/96-110.ZO.html}{Cornell LII}). The physicians declared that they would prescribe lethal medication for mentally competent, terminally ill patients to self-administer if not for Wash. Rev. Code \S~9A.36.060(1), and the plaintiffs sought a declaratory judgment that the statute violated the Fourteenth Amendment's Due Process Clause (\href{https://supreme.justia.com/cases/federal/us/521/702/}{Justia}).

\section{Procedural History}

\begin{itemize}
\item \textbf{United States District Court for the Western District of Washington (\textit{Compassion in Dying v. Washington}, 850 F. Supp. 1454 (W.D. Wash. 1994)):} Relying primarily on \textit{Planned Parenthood of Southeastern Pennsylvania v. Casey}, 505 U.S. 833 (1992), and \textit{Cruzan v. Director, Missouri Department of Health}, 497 U.S. 261 (1990), the District Court held that the Fourteenth Amendment protected a mentally competent, terminally ill adult's choice to commit physician-assisted suicide, concluded that Washington's ban imposed an ``undue burden'' on that liberty interest, and also held that the statute violated the Equal Protection Clause (\href{https://supreme.justia.com/cases/federal/us/521/702/}{Justia}).
\item \textbf{Ninth Circuit Three-Judge Panel (\textit{Compassion in Dying v. Washington}, 49 F.3d 586 (9th Cir. 1995)):} A panel of the Ninth Circuit reversed, stressing that in 205 years of constitutional history no court of final jurisdiction had ever upheld a constitutional right to aid in killing oneself (\href{https://supreme.justia.com/cases/federal/us/521/702/}{Justia}).
\item \textbf{Ninth Circuit En Banc (\textit{Compassion in Dying v. Washington}, 79 F.3d 790 (9th Cir. 1996)):} Sitting en banc, the Ninth Circuit reversed the panel and affirmed the District Court, holding that the Due Process Clause encompassed a liberty interest in ``controlling the time and manner of one's death''---a ``right to die''---and that Washington's ban was unconstitutional ``as applied to terminally ill competent adults who wish to hasten their deaths with medication prescribed by their physicians'' (\href{https://supreme.justia.com/cases/federal/us/521/702/}{Justia}). The en banc court did not reach the Equal Protection claim (\href{https://supreme.justia.com/cases/federal/us/521/702/}{Justia}).
\item \textbf{United States Supreme Court (\textit{Washington v. Glucksberg}, 521 U.S. 702 (1997)):} Certiorari was granted at 518 U.S. 1057 (1996) (\href{https://supreme.justia.com/cases/federal/us/521/702/}{Justia}). Argued January 8, 1997, alongside its New York equal-protection companion case \textit{Vacco v. Quill}, 521 U.S. 793 (1997), and decided June 26, 1997, the Supreme Court unanimously reversed and remanded (\href{https://www.oyez.org/cases/1996/96-110}{Oyez}; \href{https://www.courtlistener.com/opinion/118144/washington-v-glucksberg/}{CourtListener}).
\end{itemize}

\section{Judicial Votes}

\textbf{9--0 judgment} reversing and remanding the en banc Ninth Circuit (\href{https://supreme.justia.com/cases/federal/us/521/702/}{Justia}; \href{https://tile.loc.gov/storage-services/service/ll/usrep/usrep521/usrep521702/usrep521702.pdf}{U.S. Reports}):
\begin{itemize}
\item \textbf{Majority Opinion:} Chief Justice William H. Rehnquist delivered the opinion of the Court, joined by Justices Sandra Day O'Connor, Antonin Scalia, Anthony M. Kennedy, and Clarence Thomas (\href{https://supreme.justia.com/cases/federal/us/521/702/}{Justia}).
\item \textbf{Concurring Opinion:} Justice Sandra Day O'Connor filed a concurring opinion (applicable also to \textit{Vacco v. Quill}), which Justice Ruth Bader Ginsburg substantially endorsed in concurring in the judgments and which Justice Stephen G. Breyer joined except insofar as it joined the opinions of the Court (\href{https://supreme.justia.com/cases/federal/us/521/702/}{Justia}; \href{https://www.law.cornell.edu/supct/html/96-110.ZC4.html}{Cornell LII}).
\item \textbf{Concurrences in the Judgment(s):} Justice John Paul Stevens (concurring in the judgments), Justice David H. Souter (concurring in the judgment), Justice Ruth Bader Ginsburg (concurring in the judgments), and Justice Stephen G. Breyer (concurring in the judgments) each filed separate opinions (\href{https://supreme.justia.com/cases/federal/us/521/702/}{Justia}). There were no dissents (\href{https://supreme.justia.com/cases/federal/us/521/702/}{Justia}).
\end{itemize}

\section{Holding}

\textbf{Judgment:} Reversed and remanded (\href{https://supreme.justia.com/cases/federal/us/521/702/}{Justia}). Wash. Rev. Code \S~9A.36.060(1) does not violate the Fourteenth Amendment's Due Process Clause, either on its face or as applied to competent, terminally ill adults seeking physician-prescribed lethal medication (\href{https://supreme.justia.com/cases/federal/us/521/702/}{Justia}).

\textbf{Rule Announced:} Substantive due process analysis has two primary features: (1) the Due Process Clause specially protects only those fundamental rights and liberties that are, objectively, ``deeply rooted in this Nation's history and tradition'' and ``implicit in the concept of ordered liberty''; and (2) the Court requires a ``careful description'' of the asserted fundamental liberty interest (\href{https://supreme.justia.com/cases/federal/us/521/702/}{Justia}). If an asserted right fails this threshold test, the challenged law need only be rationally related to legitimate government interests (\href{https://supreme.justia.com/cases/federal/us/521/702/}{Justia}). Properly framed as a ``right to commit suicide which itself includes a right to assistance in doing so,'' physician-assisted suicide is not a fundamental liberty interest, and Washington's prohibition is rationally related to legitimate state interests (\href{https://supreme.justia.com/cases/federal/us/521/702/}{Justia}).

\textbf{Holding vs. Dicta:} The narrow holding is that Washington's categorical ban on assisting suicide is valid under the Fourteenth Amendment both facially and as applied to the general class of competent, terminally ill adults (\href{https://supreme.justia.com/cases/federal/us/521/702/}{Justia}). Two aspects function as important dicta or open reservations: (1) footnote 24's acknowledgment that the Court's opinion ``does not absolutely foreclose'' a more particularized as-applied challenge by an individual plaintiff, though noting such a claim would have to differ substantially from respondents' (\href{https://supreme.justia.com/cases/federal/us/521/702/}{Justia}); and (2) the majority's broad methodological assertion that the two-prong history-and-tradition test governs all substantive due process claims, which five separate concurrences either declined to endorse as an exclusive formula or expressly rejected in favor of Justice Harlan's \textit{Poe v. Ullman} framework (\href{https://supreme.justia.com/cases/federal/us/521/702/}{Justia}; \href{https://www.law.cornell.edu/supct/html/96-110.ZC4.html}{Cornell LII}).

\section{Analysis of Opinions}

\subsection{Majority Opinion (Chief Justice Rehnquist)}

Chief Justice Rehnquist began in Part I with a 700-year historical survey showing that Anglo-American common law---from Henry de Bracton and Sir William Blackstone through early colonial statutes, New York's 1828 statute, the Field Penal Code, the Model Penal Code, and modern state legislation (plus the Federal Assisted Suicide Funding Restriction Act of 1997)---had consistently condemned and prohibited assisting suicide, without any exception for those near death (\href{https://supreme.justia.com/cases/federal/us/521/702/}{Justia}). Even after states abolished harsh common-law forfeiture penalties for suicide to avoid punishing innocent families, assisting suicide remained a crime in almost every state (\href{https://supreme.justia.com/cases/federal/us/521/702/}{Justia}).

In Part II, Rehnquist articulated the Court's two-pronged substantive due process framework---requiring that an asserted fundamental right be both ``deeply rooted in this Nation's history and tradition'' (and ``implicit in the concept of ordered liberty'') and framed with a ``careful description''---to prevent judges from transforming the Due Process Clause into personal policy preferences (\href{https://supreme.justia.com/cases/federal/us/521/702/}{Justia}). Rejecting the Ninth Circuit's expansive characterizations (such as a ``right to die'' or ``liberty to shape death''), Rehnquist framed the issue narrowly as whether liberty includes ``a right to commit suicide which itself includes a right to assistance in doing so'' (\href{https://supreme.justia.com/cases/federal/us/521/702/}{Justia}). He distinguished \textit{Cruzan} as resting on the traditional common-law rule that unwanted medical touching is a battery, not abstract autonomy, and explained that \textit{Casey}'s ``mystery of human life'' passage did not convert every intimate, personal decision into a fundamental constitutional right (\href{https://supreme.justia.com/cases/federal/us/521/702/}{Justia}).

Applying rational basis review, Rehnquist concluded that Washington's ban was reasonably related to five legitimate state interests: (1) preserving human life; (2) preventing suicide as a public-health problem linked to depression and untreated pain; (3) protecting the integrity and ethics of the medical profession as healers; (4) shielding vulnerable groups---the poor, elderly, disabled, and terminally ill---from coercion, financial pressure, and societal indifference; and (5) preventing a slippery slope toward voluntary and involuntary euthanasia, citing Dutch empirical data (\href{https://supreme.justia.com/cases/federal/us/521/702/}{Justia}). He closed by affirming that the Court's ruling allowed the nation's democratic debate to continue (\href{https://supreme.justia.com/cases/federal/us/521/702/}{Justia}).

\subsection{Concurring Opinion (Justice O'Connor)}

Although joining the Court's opinion because she agreed there is no generalized right to commit suicide, Justice O'Connor wrote separately to reserve a narrower question: whether a mentally competent person experiencing great suffering has a constitutionally cognizable interest in controlling the circumstances of imminent death (\href{https://supreme.justia.com/cases/federal/us/521/702/}{Justia}). She saw no need to decide that question because both Washington and New York permitted dying patients to receive palliative medication sufficient to alleviate pain, even if it caused unconsciousness or hastened death (\href{https://supreme.justia.com/cases/federal/us/521/702/}{Justia}). Given the difficulty of defining terminal illness and ensuring voluntariness, she concluded that the democratic process---the ``laboratory'' of the States---should strike the balance (\href{https://supreme.justia.com/cases/federal/us/521/702/}{Justia}).

\subsection{Concurrence in the Judgments (Justice Stevens)}

Justice Stevens agreed that Washington's statute was not facially invalid, invoking John Donne's maxim ``No man is an island'' to recognize the State's interest in preserving community bonds and preventing ill-advised suicides (\href{https://supreme.justia.com/cases/federal/us/521/702/}{Justia}). However, analogizing to capital punishment jurisprudence, he emphasized that rejecting a facial challenge does not foreclose particularized as-applied claims (\href{https://supreme.justia.com/cases/federal/us/521/702/}{Justia}). Reading \textit{Cruzan} more broadly than the majority, Stevens argued that a terminally ill patient already on the threshold of death has a vital liberty interest in deciding \textit{how}, rather than \textit{whether}, death occurs, and in preserving dignity and surviving memories; in an individual case involving intolerable pain and a rational, uncoerced choice, the State's general interests might not outweigh that liberty interest (\href{https://supreme.justia.com/cases/federal/us/521/702/}{Justia}).

\subsection{Concurrence in the Judgment (Justice Souter)}

Justice Souter rejected the majority's rigid historical-tradition formula in favor of Justice Harlan's dissenting approach in \textit{Poe v. Ullman}, 367 U.S. 497 (1961), which treats substantive due process as a common-law process of ``reasoned judgment'' guarding against ``arbitrary impositions'' and ``purposeless restraints'' along a living tradition (\href{https://supreme.justia.com/cases/federal/us/521/702/}{Justia}). Souter recognized that the patients' asserted interest in bodily integrity and physician counsel at the end of life was serious enough to warrant heightened scrutiny (\href{https://supreme.justia.com/cases/federal/us/521/702/}{Justia}). Nevertheless, he concluded that Washington's interest in preventing involuntary suicide, abuse, and a slippery slope to euthanasia---underscored by contested empirical evidence from the Netherlands---was sufficiently weighty to defeat the claim at present, because legislatures are better equipped than courts to investigate disputed empirical facts and experiment with regulatory safeguards (\href{https://supreme.justia.com/cases/federal/us/521/702/}{Justia}).

\subsection{Concurrence in the Judgments (Justice Ginsburg)}

Justice Ginsburg issued a brief separate statement concurring in the judgments in both \textit{Glucksberg} and \textit{Vacco v. Quill} substantially for the reasons stated in Justice O'Connor's concurring opinion, thereby declining to join Chief Justice Rehnquist's majority opinion while keeping open the narrow end-of-life palliative care reservation (\href{https://supreme.justia.com/cases/federal/us/521/702/}{Justia}; \href{https://tile.loc.gov/storage-services/service/ll/usrep/usrep521/usrep521702/usrep521702.pdf}{U.S. Reports}).

\subsection{Concurrence in the Judgments (Justice Breyer)}

Justice Breyer joined Justice O'Connor's concurrence except insofar as it joined the majority opinions (\href{https://www.law.cornell.edu/supct/html/96-110.ZC4.html}{Cornell LII}). Disagreeing with the majority's narrow framing, Breyer suggested that our legal tradition might better support a ``right to die with dignity'' centered on personal control over the manner of death, medical assistance, and the avoidance of unnecessary, severe physical suffering (\href{https://www.law.cornell.edu/supct/html/96-110.ZC4.html}{Cornell LII}). He concluded, however, that the Court need not decide whether such a right is fundamental because neither Washington nor New York barred doctors from administering pain-relieving medication even at the risk of hastening death---warning that if a state law ever prevented adequate palliative care to avoid severe end-of-life pain, the Court might have to revisit its conclusion (\href{https://www.law.cornell.edu/supct/html/96-110.ZC4.html}{Cornell LII}).

\section{Examples: Future Applications}

\begin{enumerate}
\item \textbf{Same-Side (No Fundamental Right --- Access to Unapproved Experimental Drugs):} Terminally ill cancer patients sue to establish a substantive due process right to purchase and ingest experimental drugs that have passed only Phase I safety testing and lack FDA approval. Under \textit{Glucksberg}, carefully describing the asserted right as access to unapproved investigational drugs reveals no deeply rooted historical tradition protecting such access; applying rational basis review, the government's interests in protecting vulnerable patients from unsafe or ineffective therapies easily sustain the restriction (\href{https://supreme.justia.com/cases/federal/us/521/702/}{Justia}).

\item \textbf{Same-Side (No Fundamental Right --- Voluntary Active Euthanasia by a Physician):} A state criminalizes a physician's direct administration of a lethal injection at the request of a paralyzed, terminally ill patient who cannot self-administer oral medication. Under \textit{Glucksberg}, active euthanasia stands even further outside Anglo-American legal traditions than assisted suicide, and the state's interests in preserving life, maintaining medical ethics, and preventing involuntary euthanasia provide a clear rational basis for the prohibition (\href{https://supreme.justia.com/cases/federal/us/521/702/}{Justia}).

\item \textbf{Opposite-Side (Fundamental Right Protected --- Forced Life-Sustaining Medical Treatment):} A state enacts a statute prohibiting hospitals from honoring a competent, terminally ill adult's explicit refusal of a mechanical ventilator and feeding tube. Unlike assisted suicide, the right to refuse unwanted medical treatment is grounded in centuries of common-law battery and informed-consent doctrine as recognized in \textit{Cruzan} and affirmed in \textit{Glucksberg}; thus, the state's blanket override infringes a protected liberty interest (\href{https://supreme.justia.com/cases/federal/us/521/702/}{Justia}).

\item \textbf{Opposite-Side (Likely Unconstitutional --- Statutory Ban on Terminal Palliative Sedation):} A state criminalizes a physician's administration of high-dose opioid or sedative medications to relieve intractable, agonizing physical pain in a dying patient whenever the dosage carries a foreseeable risk of hastening death. Five Justices in \textit{Glucksberg} (O'Connor, Stevens, Souter, Ginsburg, and Breyer) relied expressly on the availability of palliative care---even when it hastens death---to uphold Washington's ban, and Justices O'Connor and Breyer explicitly signaled that a state law forcing a dying patient to endure severe, unavoidable physical pain would raise a grave constitutional question (\href{https://supreme.justia.com/cases/federal/us/521/702/}{Justia}; \href{https://www.law.cornell.edu/supct/html/96-110.ZC4.html}{Cornell LII}).

\item \textbf{Fence-Sitter (Level-of-Generality Clash --- Statutory Ban on Surrogacy or Embryo Destruction):} A state prohibits gestational surrogacy agreements and bars couples from discarding cryopreserved embryos created via in vitro fertilization. Under \textit{Glucksberg}'s strict ``careful description'' and history-and-tradition test, modern assisted-reproduction technologies have no centuries-old common-law pedigree, favoring rational basis review and upholding the statute; conversely, under a broader autonomy framework echoing \textit{Casey}, \textit{Skinner}, and Justice Souter's \textit{Poe} methodology, challengers would frame the liberty at a higher level of generality as fundamental procreative autonomy, triggering heightened scrutiny (\href{https://supreme.justia.com/cases/federal/us/521/702/}{Justia}; \href{https://legaljournal.princeton.edu/the-right-to-death-how-an-examination-of-washington-v-glucksberg-reveals-the-future-of-fundamental-rights/}{Princeton Legal Journal}).
\end{enumerate}

\section{Critique}

\textit{Washington v. Glucksberg} has become one of the most fiercely debated methodology precedents in constitutional law, especially after the Supreme Court invoked \textit{Glucksberg}'s ``deeply rooted in this Nation's history and tradition'' test to overrule \textit{Roe v. Wade} and \textit{Casey} in \textit{Dobbs v. Jackson Women's Health Organization}, 597 U.S. 215 (2022) (\href{https://www.supremecourt.gov/opinions/21pdf/19-1392_6j37.pdf}{Dobbs Opinion}; \href{https://legaljournal.princeton.edu/the-right-to-death-how-an-examination-of-washington-v-glucksberg-reveals-the-future-of-fundamental-rights/}{Princeton Legal Journal}).

\textbf{Progressive and Autonomy-Based Critiques:} Progressive scholars criticize both \textit{Glucksberg}'s outcome and its rigid historical methodology. Erwin Chemerinsky argued that \textit{Glucksberg} was ``tragically wrong,'' contending that the Court misapplied its own \textit{Cruzan} and \textit{Casey} precedents and failed to ask the context-specific question whether a state has a compelling interest in forcing a terminally ill patient to endure additional days or weeks of unwanted suffering when the patient already has a right to end life by refusing treatment (\href{https://repository.law.umich.edu/mlr/vol106/iss8/2/}{Michigan Law Review}). Critics further argue that demanding a ``careful description'' anchored exclusively in centuries-old common law freezes constitutional liberty in the moral assumptions of past eras---when women, racial minorities, and LGBTQ+ persons lacked political voice---and sits in deep tension with later decisions like \textit{Lawrence v. Texas} (2003) and \textit{Obergefell v. Hodges} (2015) (\href{https://repository.law.umich.edu/facarticles/259/}{Yale Kamisar, Michigan Law}; \href{https://legaljournal.princeton.edu/the-right-to-death-how-an-examination-of-washington-v-glucksberg-reveals-the-future-of-fundamental-rights/}{Princeton Legal Journal}).

\textbf{Originalist, Textualist, and Judicial-Restraint Perspectives:} Originalists and conservative scholars generally praise \textit{Glucksberg} as a vital check against unconstrained judicial policymaking under substantive due process (\href{https://fedsoc.org/case/washington-v-glucksberg}{Federalist Society}). By anchoring unenumerated rights in objective historical practice and requiring a careful description of the asserted interest, \textit{Glucksberg} restrains judges from substituting personal moral philosophy for democratic self-government---allowing states like Oregon (and later Washington itself, via voter initiative) to debate, adopt, or reject medical aid in dying through legislation (\href{https://supreme.justia.com/cases/federal/us/521/702/}{Justia}; \href{https://fedsoc.org/case/washington-v-glucksberg}{Federalist Society}). Even Yale Kamisar, who noted the tension between \textit{Glucksberg} and \textit{Lawrence}, defended judicial deference in \textit{Glucksberg} because assisted-suicide bans do not target a politically powerless or stigmatized group excluded from the democratic process (\href{https://repository.law.umich.edu/facarticles/259/}{Michigan Law}). At the same time, some legal scholars observe that \textit{Glucksberg} itself was not a purely originalist decision, since it preserved existing substantive due process precedents like \textit{Griswold}, \textit{Loving}, and \textit{Casey} while employing history and tradition as a pragmatic limiting principle (\href{https://pennlawreview.com/2025/06/10/kavanaugh-rahimi-and-the-effort-to-legitimate-originalism/}{Penn Law Review}).

\textbf{Underlying Competing Values:} The case squarely pits \textit{individual bodily autonomy, personal dignity, and relief from end-of-life suffering} against \textit{the sanctity of human life, protection of vulnerable and disabled populations from coercion, medical ethics, and democratic federalism}.

\section{Key Quotations}

\begin{quote}
``Our established method of substantive-due-process analysis has two primary features: First, we have regularly observed that the Due Process Clause specially protects those fundamental rights and liberties which are, objectively, `deeply rooted in this Nation's history and tradition,' \ldots\ and `implicit in the concept of ordered liberty,' such that `neither liberty nor justice would exist if they were sacrificed,' \textit{Palko v. Connecticut}, 302 U. S. 319, 325, 326 (1937). Second, we have required in substantive-due-process cases a `careful description' of the asserted fundamental liberty interest.''
\end{quote}
\noindent --- Chief Justice Rehnquist, Opinion of the Court, 521 U.S. at 720--721 (\href{https://supreme.justia.com/cases/federal/us/521/702/}{Justia}; \href{https://www.courtlistener.com/opinion/118144/washington-v-glucksberg/}{CourtListener}).

\begin{quote}
``That many of the rights and liberties protected by the Due Process Clause sound in personal autonomy does not warrant the sweeping conclusion that any and all important, intimate, and personal decisions are so protected, \textit{San Antonio Independent School Dist. v. Rodriguez}, 411 U. S. 1, 33--35 (1973), and \textit{Casey} did not suggest otherwise.''
\end{quote}
\noindent --- Chief Justice Rehnquist, Opinion of the Court, 521 U.S. at 727--728 (\href{https://supreme.justia.com/cases/federal/us/521/702/}{Justia}; \href{https://www.courtlistener.com/opinion/118144/washington-v-glucksberg/}{CourtListener}).

\begin{quote}
``Throughout the Nation, Americans are engaged in an earnest and profound debate about the morality, legality, and practicality of physician-assisted suicide. Our holding permits this debate to continue, as it should in a democratic society.''
\end{quote}
\noindent --- Chief Justice Rehnquist, Opinion of the Court, 521 U.S. at 735 (\href{https://supreme.justia.com/cases/federal/us/521/702/}{Justia}; \href{https://www.courtlistener.com/opinion/118144/washington-v-glucksberg/}{CourtListener}).

\begin{quote}
``There is no reason to think the democratic process will not strike the proper balance between the interests of terminally ill, mentally competent individuals who would seek to end their suffering and the State's interests in protecting those who might seek to end life mistakenly or under pressure.''
\end{quote}
\noindent --- Justice O'Connor, Concurring Opinion, 521 U.S. at 737 (\href{https://supreme.justia.com/cases/federal/us/521/702/}{Justia}; \href{https://www.courtlistener.com/opinion/118144/washington-v-glucksberg/}{CourtListener}).

\begin{quote}
``The liberty interest at stake in a case like this differs from, and is stronger than, both the common-law right to refuse medical treatment and the unbridled interest in deciding whether to live or die. It is an interest in deciding how, rather than whether, a critical threshold shall be crossed.''
\end{quote}
\noindent --- Justice Stevens, Concurring in the Judgments, 521 U.S. at 745 (\href{https://supreme.justia.com/cases/federal/us/521/702/}{Justia}; \href{https://www.courtlistener.com/opinion/118144/washington-v-glucksberg/}{CourtListener}).

\end{document}
